\documentclass[11pt,a4paper]{article}
\usepackage[margin=25mm]{geometry}
\usepackage{amsmath,amssymb,amsthm,booktabs,hyperref}
\setlength{\emergencystretch}{2em}
\newtheorem{theorem}{Theorem}\newtheorem{lemma}[theorem]{Lemma}
\newcommand{\Sym}{\operatorname{Sym}}\newcommand{\gen}[1]{\langle#1\rangle}
\title{The finite starting cases from balanced pure cuts}
\author{Human Compression project}\date{}
\begin{document}\maketitle
We acquire the proved interval of pure perfect cuts, its actual
contraction and subdivision operations, the actual-wing bound and the
exact pure-leaf criterion. A marking is fractured when every maximal
omission is intransitive. All equivalences of rooted permutation
markings below are point conjugacies taking root to root and preserving
the order of the marks. In particular degree-six wings are not
identified under the exceptional outer automorphism of $S_6$.

Call a rooted marking of degree $2p$ and rank $p$ a \emph{wing} if it is
transitive, fractured and C, its root is moved by its first mark and fixed
by all later marks, and adjoining a pure transposition from the root to
a new point gives a C-marking. Write $w_p$ for the number of its rooted
classes. The orientation of a wing starts at its root; the left wing in
a joined marking is therefore reversed.

\begin{theorem}\label{balanced}
For $\kappa\ge1$, the rank-$\kappa+3$ C-markings of
$\Sym(2\kappa+3)$, up to marked-group isomorphism and duality,
are in bijection with unordered pairs of wings of ranks $p,q\ge1$
satisfying $p+q=\kappa+1$. In particular their number is
\[
 c_\kappa=\sum_{\substack{p<q\\p+q=\kappa+1}}w_pw_q+
 \begin{cases}w_h(w_h+1)/2,&\kappa+1=2h,\\0,&\kappa+1\text{ odd}.
 \end{cases}                                                   \tag{1}
\]
\end{theorem}
\begin{proof}
Put $n=2r-3$. For $n\ge7$ the acquired high-rank theorem makes every
natural symmetric marking fractured, and the acquired pure-cut theorem
gives an internal pure perfect cut. All pure perfect cuts form one
interval. At a cut $i$ with prefix support $A_i$ put
$d(i)=2i-|A_i|$. Its suffix surplus is $1-d(i)$. Consecutive pure
cuts add one point and one mark to the prefix, so
$d(i+1)=d(i)+1$.

At the first pure cut we have $d\le0$. Otherwise its prefix is
symmetric by the actual-wing bound, and its old pure-leaf extension
satisfies the full kernel tests. The acquired backward propagation
lemma makes the preceding mark pure. Deleting the pendant endpoint
leaves the preceding prefix transitive, so this is a pure perfect cut
of the whole marking, a contradiction. A positive surplus requires
prefix rank at least two, so the preceding label exists. Reversal
shows that the suffix surplus at the last pure cut is nonpositive;
there $d\ge1$. There is exactly one consecutive pair with surpluses
$0,1$.

Its adjacent geometry consists of disjoint supports $A,\{c\},B$,
preceding and following groups $L,T$, and marks $(a\ c),(c\ b)$.
If $L,T$ have ranks $p,q$, the two surplus equations give
$|A|=2p$, $|B|=2q$, and $p+q=r-2$. Both ranks are positive,
so both cuts are internal. Orienting $L$ from its last mark gives
its root $a$; orient $T$ from its first mark with root $b$.
Distant commutation fixes each root under every later wing mark.
Each wing is C by the original IP, and its pure-leaf extension is
an old interval. Each is fractured: an omission transitive on its
wing would remain transitive on the entire graph through the two
bridges and the other wing, contradicting fracture of the parent.
Thus these are wings in the stated sense.

Conversely take any two wings on disjoint supports $A,B$ rooted at
$a,b$. Join them first by $q=(a\ b)$. Their pure-leaf extensions
ensure that intervals ending or beginning at $q$ are C. For intervals
crossing $q$, the actual one-bridge intersection rectangle from the
acquired contraction proof and induction on interval length give IP.
For clarity that rectangle is
\[
 \gen{P_0,Q_1,q}\cap\gen{P_1,Q_0,q}=\gen{P_0,Q_0,q}
       \quad(P_0\le P_1\le\Sym(A),\ Q_0\le Q_1\le\Sym(B)).
\]
The full transitive wings generate $\Sym(A\cup B)$ with $q$.
The mark $q$ is a pure perfect cut; its acquired subdivision gives
$(a\ c),(c\ b)$ at a new point and preserves IP. The resulting
marking has degree $2p+2q+1$, rank $p+q+2$, and its two indicated
pure cuts have surpluses $0,1$. They are therefore precisely its
unique pair just described, even if the pure interval is longer.
These two constructions are inverse on ordered rooted wing classes.

An ordered point conjugacy must preserve the unique balanced pair;
it restricts to rooted conjugacies of its two wings. Conversely these
rooted conjugacies extend by sending the middle point to the middle
point. Reversal exchanges the two wings. For degree at least seven,
marked-group isomorphisms of natural symmetric groups are precisely
point conjugacies, by the acquired automorphism argument. In degree
five the acquired simplex class is obtained from the unique degree-two,
rank-one wing twice. This handles $\kappa=1$ separately and proves
(1) in every case.
\end{proof}

\begin{lemma}\label{search}
The following finite search lists each wing of a fixed even degree $n$
exactly once up to rooted point conjugacy.
\end{lemma}
\begin{proof}
Set $p=n/2$ and normalize the root to $1$. List every nonidentity
involution of $\Sym(n)$ by conjugating the products of $j$ disjoint
transpositions, for $1\le j\le n/2$. The first mark moves $1$;
all subsequent marks fix $1$. At each prefix try involutions outside
the group already generated, commuting with every earlier mark except
the last, and not commuting with that last mark.

This last requirement loses no transitive anchored string marking.
The component of its noncommutation diagram containing the first mark
is transitive: every other component fixes the root and commutes with
that component, and therefore fixes its root orbit pointwise. The full
orbit is consequently that component's orbit. Every other component
would then consist of identity permutations, which is impossible.
Thus the noncommutation diagram is the whole path.

Use one orbit representative under the point stabilizer of $1$ at the
first step, and thereafter under its subgroup centralizing the entire
chosen prefix. This gives exactly the rooted conjugacy classes of
prefixes, recursively. Every condition imposed is invariant under
these conjugations. Test IP of each new interval by the standard end
criterion, using the previously verified intervals. Thus all prefixes
retained are C. A separate terminal oracle checks all pairs of all
subsets on each retained transitive marking.

A fractured final marking must have every prefix mark essential to
the prefix orbit partition. If deleting that mark left the partition
unchanged, adding the remaining marks would merge the same blocks
in either case and leave the final omission transitive. Reject exactly
such inessential prefixes. If a prefix generates $H$ and has length
$l$, also reject when $[\Sym(n):H]<2^{p-l}$: independence of each
remaining mark would multiply the order by at least two.

For every enlarged prefix $P_j$, compute $X_{j-1}=1^{P_{j-1}}$
and impose the actual subgroup inclusion
\[
                 P_j\cap\Sym(X_{j-1})\le P_{j-1}.       \tag{2}
\]
These are exactly the acquired full leaf-kernel tests, not tests on
an unproved quotient action. Finally retain only transitive terminal
markings. Prefix IP, orbit-essentiality and (2) then give respectively
C, fracture and pure-leaf validity. Every possible wing survives every
necessary test, while rooted orbit representatives eliminate precisely
the declared equivalence. For efficiency the candidate list at a prefix retains just the
root-fixing involutions commuting with all marks except the last.
At length one it is the full root-fixing list; the child list is
obtained by imposing commutation with the old last mark. This
inductive invariant reproduces the literal distant-commutation test.
This proves completeness and uniqueness.
\end{proof}

\begin{theorem}\label{finite}
The first six wing counts are
\[
                    (w_1,\ldots,w_6)=(1,1,6,3,11,19).
\]
Consequently the first six stable constants in the acquired counting
theorem are $(1,1,7,9,35,48)$.
\end{theorem}
\begin{proof}
The complete executions of Lemma~\ref{search} give the following table.
The prefix counts include the empty prefix. An intransitive terminal
prefix is not a wing and contributes only to that column.
\begin{center}\small
\begin{tabular}{crrcl}
\toprule
$n$ & involutions & root-fixing & prefixes by length & wings\\
\midrule
2&1&0&$1,1$&1\\
4&9&3&$1,2,2$&1\\
6&75&25&$1,3,12,11$&6\\
8&763&231&$1,4,53,95,29$&3\\
10&9495&2619&$1,5,187,1054,415,156$&11\\
12&140151&35695&$1,6,552,10218,7912,3149,617$&19\\
\bottomrule
\end{tabular}\end{center}
Every run terminated with its unique completion marker and empty error
output. The largest completed run took about 91 seconds. Two earlier
180-second caps supply no part of the count; their outputs are retained.
The final candidate-list cache preserves the literal distant-commutation
predicate and reproduced the entire degree-ten output byte for byte.
Moving the terminal transitivity test before the separate all-subset IP
oracle only discards intransitive groups outside the target; every
retained wing still receives that full oracle. A further proposed
necessary-prefix test made no additional rejections on its control and
is not used in the chosen final engine.

The full literal representatives are retained. Independently, a labelled
degree-six search tries every root-moving first involution and every
root-fixing second and third involution, imposing distant commutation,
then checks actual transitivity, all maximal omissions, every subset
intersection and each full leaf kernel. It uses neither rooted orbit
pruning, the prefix orbit test, the order bound nor the incremental IP
test. Its surviving rooted classes agree exactly with the main six.
A separate permutation-array implementation checks every subgroup closure
and leaf kernel in degrees two, four and six.

A rooted connected marking has a canonical literal array obtained by
breadth-first traversal from its root, trying the labels in their fixed
order. Two such arrays coincide exactly when a rooted point conjugacy
exists. This independently verifies distinctness of every wing list.
For each unordered pair in Theorem~\ref{balanced}, construct its actual
joined tuple. Minimizing the analogous code over every point and the
two orientations verifies that all resulting boundary tuples are distinct
modulo point conjugacy and reversal. There are respectively
\[
 1,\quad 1,\quad 6+1=7,\quad 3+6=9,\quad
 11+3+21=35,\quad 19+11+18=48.
\]
Their generation and IP follow from the proved gluing construction;
separate finite controls check these actual tuples as well. Completeness
is Theorem~\ref{balanced} together with Lemma~\ref{search}, not agreement
with the historical sequence. The acquired stability theorem transfers
these boundary counts to every degree in the stated stable ranges.
\end{proof}

\paragraph{Attribution and scope.}
The program is an explicit derivative of the previously reviewed
degree-ten anchored search, with full rather than even leaf kernels,
all first involutions, and all transitive terminal groups. The structural
balanced-pair argument and its use to compute the finite bases are new
to this recovery, without an assertion of historical priority. The
original stable counts remain attributed to Cameron, Fernandes and
Leemans and their preceding collaborators. No external atlas completeness
claim is used. All actual source inputs, completed runs, capped attempts
and changes of engine are pinned separately. This component evaluates
the original six constants; it does not begin a further sequence search,
new conjecture programme, Lean development or compression measurement.
\end{document}
